% --- Archon physics formalization source begin ---
% archon:physics
% archon:covers PhyXMiniProblems/problem_phyx_mini_0736.lean
% archon:source-report reports/phyx_mini/problem_phyx_mini_0736.source.json

\chapter{Physics problem phyx\_mini\_0736}
\label{ch:PhyXMiniProblems_problem_phyx_mini_0736}

\paragraph{Problem source.}
\#\# Physical scenario

A wheel with a radius of \$45.0\textbackslash{} cm\$ rolls without slipping along a horizontal floor. At time \$t\_1\$, the dot \$P\$ painted on the rim of the wheel is at the point of contact between the wheel and the floor. At a later time \$t\_2\$, the wheel has rolled through one-half of a revolution.

\#\# Question

What is the angle (relative to the floor) of the displacement of \$P\$?

\#\# Answer choices

- A: A: \$25.5\textasciicircum{}\textbackslash{}circ\$

- B: B: \$27.5\textasciicircum{}\textbackslash{}circ\$

- C: C: \$32.5\textasciicircum{}\textbackslash{}circ\$

- D: D: \$34.5\textasciicircum{}\textbackslash{}circ\$

\#\# Figure caption (auxiliary; use the image as primary evidence)

The image features two diagrams of a wheel at different times, labeled "At time \textbackslash{}(t\_1\textbackslash{})" and "At time \textbackslash{}(t\_2\textbackslash{})." 

1. **At time \textbackslash{}(t\_1\textbackslash{}):**
   - A wheel is in contact with a textured surface at the bottom.
   - Point \textbackslash{}(P\textbackslash{}) is indicated at the bottom of the wheel touching the surface.

2. **At time \textbackslash{}(t\_2\textbackslash{}):**
   - The same wheel is shown, but it has rotated so that point \textbackslash{}(P\textbackslash{}) is now at the top of the wheel, no longer touching the surface.

Both diagrams illustrate the rotation of the wheel over time, specifically focusing on the change in position of point \textbackslash{}(P\textbackslash{}) as the wheel rolls along the surface. The textured surface is consistent in both diagrams.

\#\# Dataset metadata

Rotational Motion; Spatial Relation Reasoning, Physical Model Grounding Reasoning

\paragraph{Recorded answer/context.}
C: C: \$32.5\textasciicircum{}\textbackslash{}circ\$

\paragraph{Figure/image path.}
/root/proposal\_for\_physic/science-mango/phyx\_mini\_run/phyx\_data/test\_image/736.png

\paragraph{Formalization target.}
create a compiling Lean file with sorry bodies at `PhyXMiniProblems/problem\_phyx\_mini\_0736.lean`. The Lean declarations must preserve the physical quantities, dimensions or dimensional roles, figure labels, governing-law hypotheses, and final relation expressed by this problem.
Use Mathlib/Physlib names found through LeanExplore where available. If a physics API is missing, introduce faithful local abstractions rather than scalar placeholder aliases.

\begin{theorem}[Physics formalization target]
\label{thm:physics:phyx_mini_0736:target}
\lean{PhyXMiniProblems.ProblemPhyXMini0736.problem_phyx_mini_0736}
\uses{def:physics:phyx-mini-0736:phyxminiproblems-problemphyxmini0736-rollingwheelsetup, def:physics:phyx-mini-0736:phyxminiproblems-problemphyxmini0736-markedpointdisplacementangleradians, def:physics:phyx-mini-0736:phyxminiproblems-problemphyxmini0736-matchesproblemstatement, def:physics:phyx-mini-0736:phyxminiproblems-problemphyxmini0736-matchesprimaryfigure, def:physics:phyx-mini-0736:phyxminiproblems-problemphyxmini0736-hasphysicalwheelparameters, def:physics:phyx-mini-0736:phyxminiproblems-problemphyxmini0736-satisfiesrigidwheelgeometry, def:physics:phyx-mini-0736:phyxminiproblems-problemphyxmini0736-satisfiesrollingwithoutslipping, def:physics:phyx-mini-0736:phyxminiproblems-problemphyxmini0736-recordedanswerchoice, def:physics:phyx-mini-0736:phyxminiproblems-problemphyxmini0736-matchesdisplayedprecision, def:physics:phyx-mini-0736:phyxminiproblems-problemphyxmini0736-isuniqueclosestanswer}
The assigned autoformalize agent should translate this physics problem into Lean declarations in the covered file, with theorem and lemma proof bodies written as `by sorry`.
\end{theorem}
\begin{proof}
This is an autoformalization task, not a proof task. Produce statements that can later be proved by the physics prover without weakening the physical model.
\end{proof}

% --- Archon declaration topology begin ---
\section{Declaration topology}
\begin{definition}[Length Quantity]
  \label{def:physics:phyx-mini-0736:phyxminiproblems-problemphyxmini0736-lengthquantity}
  \lean{PhyXMiniProblems.ProblemPhyXMini0736.LengthQuantity}
  A nonnegative physical length, used for the wheel radius
\end{definition}
\begin{proof}
  This is the declared model definition of Length Quantity.
\end{proof}
\begin{definition}[Signed Length Quantity]
  \label{def:physics:phyx-mini-0736:phyxminiproblems-problemphyxmini0736-signedlengthquantity}
  \lean{PhyXMiniProblems.ProblemPhyXMini0736.SignedLengthQuantity}
  A signed physical length, used for a coordinate relative to a chosen origin
\end{definition}
\begin{proof}
  This is the declared model definition of Signed Length Quantity.
\end{proof}
\begin{definition}[length Readout]
  \label{def:physics:phyx-mini-0736:phyxminiproblems-problemphyxmini0736-lengthreadout}
  \lean{PhyXMiniProblems.ProblemPhyXMini0736.lengthReadout}
  \uses{def:physics:phyx-mini-0736:phyxminiproblems-problemphyxmini0736-lengthquantity}
  Scalar readout of a nonnegative physical length in a coherent unit system
\end{definition}
\begin{proof}
  This is the declared model definition of length Readout.
\end{proof}
\begin{definition}[signed Length Readout]
  \label{def:physics:phyx-mini-0736:phyxminiproblems-problemphyxmini0736-signedlengthreadout}
  \lean{PhyXMiniProblems.ProblemPhyXMini0736.signedLengthReadout}
  \uses{def:physics:phyx-mini-0736:phyxminiproblems-problemphyxmini0736-signedlengthquantity}
  Scalar readout of a signed physical coordinate in a coherent unit system
\end{definition}
\begin{proof}
  This is the declared model definition of signed Length Readout.
\end{proof}
\begin{definition}[centimeter Unit Choices]
  \label{def:physics:phyx-mini-0736:phyxminiproblems-problemphyxmini0736-centimeterunitchoices}
  \lean{PhyXMiniProblems.ProblemPhyXMini0736.centimeterUnitChoices}
  SI base units with centimeters selected as the length unit
\end{definition}
\begin{proof}
  This is the declared model definition of centimeter Unit Choices.
\end{proof}
\begin{definition}[length In Centimeters]
  \label{def:physics:phyx-mini-0736:phyxminiproblems-problemphyxmini0736-lengthincentimeters}
  \lean{PhyXMiniProblems.ProblemPhyXMini0736.lengthInCentimeters}
  \uses{def:physics:phyx-mini-0736:phyxminiproblems-problemphyxmini0736-lengthquantity, def:physics:phyx-mini-0736:phyxminiproblems-problemphyxmini0736-lengthreadout, def:physics:phyx-mini-0736:phyxminiproblems-problemphyxmini0736-centimeterunitchoices}
  Read a physical radius in centimeters
\end{definition}
\begin{proof}
  This is the declared model definition of length In Centimeters.
\end{proof}
\begin{definition}[Plane Axis]
  \label{def:physics:phyx-mini-0736:phyxminiproblems-problemphyxmini0736-planeaxis}
  \lean{PhyXMiniProblems.ProblemPhyXMini0736.PlaneAxis}
  The two coordinate directions in the plane of the wheel
\end{definition}
\begin{proof}
  This is the declared model definition of Plane Axis.
\end{proof}
\begin{definition}[Plane Position]
  \label{def:physics:phyx-mini-0736:phyxminiproblems-problemphyxmini0736-planeposition}
  \lean{PhyXMiniProblems.ProblemPhyXMini0736.PlanePosition}
  \uses{def:physics:phyx-mini-0736:phyxminiproblems-problemphyxmini0736-signedlengthquantity, def:physics:phyx-mini-0736:phyxminiproblems-problemphyxmini0736-planeaxis}
  A planar physical position, represented by one signed length per axis
\end{definition}
\begin{proof}
  This is the declared model definition of Plane Position.
\end{proof}
\begin{definition}[Wheel Instant]
  \label{def:physics:phyx-mini-0736:phyxminiproblems-problemphyxmini0736-wheelinstant}
  \lean{PhyXMiniProblems.ProblemPhyXMini0736.WheelInstant}
  The two observation times named in the figure
\end{definition}
\begin{proof}
  This is the declared model definition of Wheel Instant.
\end{proof}
\begin{definition}[Figure Time Label]
  \label{def:physics:phyx-mini-0736:phyxminiproblems-problemphyxmini0736-figuretimelabel}
  \lean{PhyXMiniProblems.ProblemPhyXMini0736.FigureTimeLabel}
  Literal time labels printed beneath the two panels
\end{definition}
\begin{proof}
  This is the declared model definition of Figure Time Label.
\end{proof}
\begin{definition}[Marked Point Label]
  \label{def:physics:phyx-mini-0736:phyxminiproblems-problemphyxmini0736-markedpointlabel}
  \lean{PhyXMiniProblems.ProblemPhyXMini0736.MarkedPointLabel}
  The label attached to the painted rim point
\end{definition}
\begin{proof}
  This is the declared model definition of Marked Point Label.
\end{proof}
\begin{definition}[Rim Location]
  \label{def:physics:phyx-mini-0736:phyxminiproblems-problemphyxmini0736-rimlocation}
  \lean{PhyXMiniProblems.ProblemPhyXMini0736.RimLocation}
  Qualitative locations of the marked point in the two supplied panels
\end{definition}
\begin{proof}
  This is the declared model definition of Rim Location.
\end{proof}
\begin{definition}[Floor Orientation]
  \label{def:physics:phyx-mini-0736:phyxminiproblems-problemphyxmini0736-floororientation}
  \lean{PhyXMiniProblems.ProblemPhyXMini0736.FloorOrientation}
  Orientation of the supporting floor
\end{definition}
\begin{proof}
  This is the declared model definition of Floor Orientation.
\end{proof}
\begin{definition}[Rolling Direction]
  \label{def:physics:phyx-mini-0736:phyxminiproblems-problemphyxmini0736-rollingdirection}
  \lean{PhyXMiniProblems.ProblemPhyXMini0736.RollingDirection}
  Horizontal rolling direction indicated by the left-to-right panel layout
\end{definition}
\begin{proof}
  This is the declared model definition of Rolling Direction.
\end{proof}
\begin{definition}[Requested Observable]
  \label{def:physics:phyx-mini-0736:phyxminiproblems-problemphyxmini0736-requestedobservable}
  \lean{PhyXMiniProblems.ProblemPhyXMini0736.RequestedObservable}
  The physical observable requested in the question
\end{definition}
\begin{proof}
  This is the declared model definition of Requested Observable.
\end{proof}
\begin{definition}[Rolling Wheel Figure]
  \label{def:physics:phyx-mini-0736:phyxminiproblems-problemphyxmini0736-rollingwheelfigure}
  \lean{PhyXMiniProblems.ProblemPhyXMini0736.RollingWheelFigure}
  \uses{def:physics:phyx-mini-0736:phyxminiproblems-problemphyxmini0736-wheelinstant, def:physics:phyx-mini-0736:phyxminiproblems-problemphyxmini0736-figuretimelabel, def:physics:phyx-mini-0736:phyxminiproblems-problemphyxmini0736-markedpointlabel, def:physics:phyx-mini-0736:phyxminiproblems-problemphyxmini0736-rimlocation}
  Qualitative information transcribed from the two-panel primary image
\end{definition}
\begin{proof}
  This is the declared model definition of Rolling Wheel Figure.
\end{proof}
\begin{definition}[Rolling Wheel Setup]
  \label{def:physics:phyx-mini-0736:phyxminiproblems-problemphyxmini0736-rollingwheelsetup}
  \lean{PhyXMiniProblems.ProblemPhyXMini0736.RollingWheelSetup}
  \uses{def:physics:phyx-mini-0736:phyxminiproblems-problemphyxmini0736-lengthquantity, def:physics:phyx-mini-0736:phyxminiproblems-problemphyxmini0736-planeposition, def:physics:phyx-mini-0736:phyxminiproblems-problemphyxmini0736-wheelinstant, def:physics:phyx-mini-0736:phyxminiproblems-problemphyxmini0736-floororientation, def:physics:phyx-mini-0736:phyxminiproblems-problemphyxmini0736-rollingdirection, def:physics:phyx-mini-0736:phyxminiproblems-problemphyxmini0736-requestedobservable, def:physics:phyx-mini-0736:phyxminiproblems-problemphyxmini0736-rollingwheelfigure}
  The independent physical quantities and trajectories in the scenario. Neither the displacement direction nor any answer choice is stored here
\end{definition}
\begin{proof}
  This is the declared model definition of Rolling Wheel Setup.
\end{proof}
\begin{definition}[coordinate Readout]
  \label{def:physics:phyx-mini-0736:phyxminiproblems-problemphyxmini0736-coordinatereadout}
  \lean{PhyXMiniProblems.ProblemPhyXMini0736.coordinateReadout}
  \uses{def:physics:phyx-mini-0736:phyxminiproblems-problemphyxmini0736-signedlengthreadout, def:physics:phyx-mini-0736:phyxminiproblems-problemphyxmini0736-planeaxis, def:physics:phyx-mini-0736:phyxminiproblems-problemphyxmini0736-planeposition}
  Read one coordinate of a dimensionful planar position
\end{definition}
\begin{proof}
  This is the declared model definition of coordinate Readout.
\end{proof}
\begin{definition}[displacement Coordinate Readout]
  \label{def:physics:phyx-mini-0736:phyxminiproblems-problemphyxmini0736-displacementcoordinatereadout}
  \lean{PhyXMiniProblems.ProblemPhyXMini0736.displacementCoordinateReadout}
  \uses{def:physics:phyx-mini-0736:phyxminiproblems-problemphyxmini0736-planeaxis, def:physics:phyx-mini-0736:phyxminiproblems-problemphyxmini0736-planeposition, def:physics:phyx-mini-0736:phyxminiproblems-problemphyxmini0736-coordinatereadout}
  Coordinate displacement between two dimensionful planar positions
\end{definition}
\begin{proof}
  This is the declared model definition of displacement Coordinate Readout.
\end{proof}
\begin{definition}[center Displacement Readout]
  \label{def:physics:phyx-mini-0736:phyxminiproblems-problemphyxmini0736-centerdisplacementreadout}
  \lean{PhyXMiniProblems.ProblemPhyXMini0736.centerDisplacementReadout}
  \uses{def:physics:phyx-mini-0736:phyxminiproblems-problemphyxmini0736-planeaxis, def:physics:phyx-mini-0736:phyxminiproblems-problemphyxmini0736-rollingwheelsetup, def:physics:phyx-mini-0736:phyxminiproblems-problemphyxmini0736-displacementcoordinatereadout}
  Displacement of the wheel center from t₁ to t₂
\end{definition}
\begin{proof}
  This is the declared model definition of center Displacement Readout.
\end{proof}
\begin{definition}[marked Point Displacement Readout]
  \label{def:physics:phyx-mini-0736:phyxminiproblems-problemphyxmini0736-markedpointdisplacementreadout}
  \lean{PhyXMiniProblems.ProblemPhyXMini0736.markedPointDisplacementReadout}
  \uses{def:physics:phyx-mini-0736:phyxminiproblems-problemphyxmini0736-planeaxis, def:physics:phyx-mini-0736:phyxminiproblems-problemphyxmini0736-rollingwheelsetup, def:physics:phyx-mini-0736:phyxminiproblems-problemphyxmini0736-displacementcoordinatereadout}
  Displacement of the painted point P from t₁ to t₂
\end{definition}
\begin{proof}
  This is the declared model definition of marked Point Displacement Readout.
\end{proof}
\begin{definition}[marked Point Displacement Angle Radians]
  \label{def:physics:phyx-mini-0736:phyxminiproblems-problemphyxmini0736-markedpointdisplacementangleradians}
  \lean{PhyXMiniProblems.ProblemPhyXMini0736.markedPointDisplacementAngleRadians}
  \uses{def:physics:phyx-mini-0736:phyxminiproblems-problemphyxmini0736-rollingwheelsetup, def:physics:phyx-mini-0736:phyxminiproblems-problemphyxmini0736-markedpointdisplacementreadout}
  Angle of the marked point's displacement above the positive horizontal floor direction. The physical assumptions below imply a positive horizontal component, selecting the ordinary Real.arctan branch
\end{definition}
\begin{proof}
  This is the declared model definition of marked Point Displacement Angle Radians.
\end{proof}
\begin{definition}[radians To Degrees]
  \label{def:physics:phyx-mini-0736:phyxminiproblems-problemphyxmini0736-radianstodegrees}
  \lean{PhyXMiniProblems.ProblemPhyXMini0736.radiansToDegrees}
  Convert a dimensionless radian readout to degrees
\end{definition}
\begin{proof}
  This is the declared model definition of radians To Degrees.
\end{proof}
\begin{definition}[Matches Problem Statement]
  \label{def:physics:phyx-mini-0736:phyxminiproblems-problemphyxmini0736-matchesproblemstatement}
  \lean{PhyXMiniProblems.ProblemPhyXMini0736.MatchesProblemStatement}
  \uses{def:physics:phyx-mini-0736:phyxminiproblems-problemphyxmini0736-lengthincentimeters, def:physics:phyx-mini-0736:phyxminiproblems-problemphyxmini0736-rollingwheelsetup}
  Problem-text data: a 45.0 cm wheel on a horizontal floor completes a half-revolution, and the question asks for the direction of P's displacement. No displacement component or answer angle is included
\end{definition}
\begin{proof}
  This is the declared model definition of Matches Problem Statement.
\end{proof}
\begin{definition}[Matches Primary Figure]
  \label{def:physics:phyx-mini-0736:phyxminiproblems-problemphyxmini0736-matchesprimaryfigure}
  \lean{PhyXMiniProblems.ProblemPhyXMini0736.MatchesPrimaryFigure}
  \uses{def:physics:phyx-mini-0736:phyxminiproblems-problemphyxmini0736-rollingwheelsetup}
  Primary-image evidence: two labeled panels show the same wheel and floor; P is the bottom contact point at t₁, the topmost rim point at t₂, and the later panel lies to the right. These qualitative readouts contain no answer angle or numerical displacement
\end{definition}
\begin{proof}
  This is the declared model definition of Matches Primary Figure.
\end{proof}
\begin{definition}[Has Physical Wheel Parameters]
  \label{def:physics:phyx-mini-0736:phyxminiproblems-problemphyxmini0736-hasphysicalwheelparameters}
  \lean{PhyXMiniProblems.ProblemPhyXMini0736.HasPhysicalWheelParameters}
  \uses{def:physics:phyx-mini-0736:phyxminiproblems-problemphyxmini0736-lengthreadout, def:physics:phyx-mini-0736:phyxminiproblems-problemphyxmini0736-rollingwheelsetup}
  Positivity needed to cancel the physical radius and select the angle branch
\end{definition}
\begin{proof}
  This is the declared model definition of Has Physical Wheel Parameters.
\end{proof}
\begin{definition}[Satisfies Rigid Wheel Geometry]
  \label{def:physics:phyx-mini-0736:phyxminiproblems-problemphyxmini0736-satisfiesrigidwheelgeometry}
  \lean{PhyXMiniProblems.ProblemPhyXMini0736.SatisfiesRigidWheelGeometry}
  \uses{def:physics:phyx-mini-0736:phyxminiproblems-problemphyxmini0736-lengthreadout, def:physics:phyx-mini-0736:phyxminiproblems-problemphyxmini0736-rollingwheelsetup, def:physics:phyx-mini-0736:phyxminiproblems-problemphyxmini0736-coordinatereadout}
  Rigid-wheel geometry relative to the floor coordinate frame. The center stays one radius above the floor. A qualitative bottom/top rim location places P one radius below/above the center, respectively. These laws are stated in every coherent unit system and do not mention the requested angle
\end{definition}
\begin{proof}
  This is the declared model definition of Satisfies Rigid Wheel Geometry.
\end{proof}
\begin{definition}[Satisfies Rolling Without Slipping]
  \label{def:physics:phyx-mini-0736:phyxminiproblems-problemphyxmini0736-satisfiesrollingwithoutslipping}
  \lean{PhyXMiniProblems.ProblemPhyXMini0736.SatisfiesRollingWithoutSlipping}
  \uses{def:physics:phyx-mini-0736:phyxminiproblems-problemphyxmini0736-lengthreadout, def:physics:phyx-mini-0736:phyxminiproblems-problemphyxmini0736-rollingwheelsetup, def:physics:phyx-mini-0736:phyxminiproblems-problemphyxmini0736-centerdisplacementreadout}
  Rolling without slipping: for the rightward motion shown, the center's horizontal travel is the wheel radius times the positive rotation magnitude. This is the governing kinematic law, not a statement of P's displacement or its angle
\end{definition}
\begin{proof}
  This is the declared model definition of Satisfies Rolling Without Slipping.
\end{proof}
\begin{definition}[Answer Choice]
  \label{def:physics:phyx-mini-0736:phyxminiproblems-problemphyxmini0736-answerchoice}
  \lean{PhyXMiniProblems.ProblemPhyXMini0736.AnswerChoice}
  Labels of the four displayed answers
\end{definition}
\begin{proof}
  This is the declared model definition of Answer Choice.
\end{proof}
\begin{definition}[degrees]
  \label{def:physics:phyx-mini-0736:phyxminiproblems-problemphyxmini0736-answerchoice-degrees}
  \lean{PhyXMiniProblems.ProblemPhyXMini0736.AnswerChoice.degrees}
  \uses{def:physics:phyx-mini-0736:phyxminiproblems-problemphyxmini0736-answerchoice}
  Angle in degrees printed beside each answer label
\end{definition}
\begin{proof}
  This is the declared model definition of degrees.
\end{proof}
\begin{definition}[recorded Answer Choice]
  \label{def:physics:phyx-mini-0736:phyxminiproblems-problemphyxmini0736-recordedanswerchoice}
  \lean{PhyXMiniProblems.ProblemPhyXMini0736.recordedAnswerChoice}
  \uses{def:physics:phyx-mini-0736:phyxminiproblems-problemphyxmini0736-answerchoice}
  Dataset metadata: the recorded answer is choice C; this is never a premise
\end{definition}
\begin{proof}
  This is the declared model definition of recorded Answer Choice.
\end{proof}
\begin{definition}[Matches Displayed Precision]
  \label{def:physics:phyx-mini-0736:phyxminiproblems-problemphyxmini0736-matchesdisplayedprecision}
  \lean{PhyXMiniProblems.ProblemPhyXMini0736.MatchesDisplayedPrecision}
  \uses{def:physics:phyx-mini-0736:phyxminiproblems-problemphyxmini0736-radianstodegrees, def:physics:phyx-mini-0736:phyxminiproblems-problemphyxmini0736-answerchoice}
  Agreement with a degree value displayed to the nearest tenth. The half-tenth-degree tolerance models the precision of the answer list
\end{definition}
\begin{proof}
  This is the declared model definition of Matches Displayed Precision.
\end{proof}
\begin{definition}[Is Unique Closest Answer]
  \label{def:physics:phyx-mini-0736:phyxminiproblems-problemphyxmini0736-isuniqueclosestanswer}
  \lean{PhyXMiniProblems.ProblemPhyXMini0736.IsUniqueClosestAnswer}
  \uses{def:physics:phyx-mini-0736:phyxminiproblems-problemphyxmini0736-radianstodegrees, def:physics:phyx-mini-0736:phyxminiproblems-problemphyxmini0736-answerchoice}
  A displayed answer is uniquely closest to the physical angle
\end{definition}
\begin{proof}
  This is the declared model definition of Is Unique Closest Answer.
\end{proof}
\begin{lemma}[marked Point displacement components]
  \label{lem:physics:phyx-mini-0736:phyxminiproblems-problemphyxmini0736-markedpoint-displacement-components}
  \lean{PhyXMiniProblems.ProblemPhyXMini0736.markedPoint_displacement_components}
  \uses{def:physics:phyx-mini-0736:phyxminiproblems-problemphyxmini0736-lengthreadout, def:physics:phyx-mini-0736:phyxminiproblems-problemphyxmini0736-rollingwheelsetup, def:physics:phyx-mini-0736:phyxminiproblems-problemphyxmini0736-markedpointdisplacementreadout, def:physics:phyx-mini-0736:phyxminiproblems-problemphyxmini0736-matchesproblemstatement, def:physics:phyx-mini-0736:phyxminiproblems-problemphyxmini0736-matchesprimaryfigure, def:physics:phyx-mini-0736:phyxminiproblems-problemphyxmini0736-satisfiesrigidwheelgeometry, def:physics:phyx-mini-0736:phyxminiproblems-problemphyxmini0736-satisfiesrollingwithoutslipping}
  The no-slip and rigid-wheel relations imply that, in any coherent length unit, the displacement of P is (π R, 2 R). This is a derived kinematic result rather than a premise
\end{lemma}
\begin{proof}
  The conclusion follows from the governing relations and figure readouts named in the statement of marked Point displacement components.
\end{proof}
% --- Archon declaration topology end ---
% --- Archon physics formalization source end ---
